\documentclass[a4paper,12pt]{article}
\usepackage[margin=2.5cm]{geometry}
\usepackage{graphicx}
\usepackage{booktabs}
\usepackage{amsmath}
\usepackage{amssymb}
\usepackage{caption}
\usepackage{subcaption}
\usepackage{hyperref}
\usepackage{xcolor}
\usepackage{tabularx}

\title{\textbf{Result 12: Mobility Impact on Average End-to-End Delay}\\
\large ACLIF --- Autonomous Cross-Layer Intelligence Framework}
\author{B.~Tharuni~Sri~Sai\\
\small School of Computer Science and Engineering, VIT-AP University, Amaravati, India\\
\small ORCID: 0009-0004-9561-8985}
\date{\today}

\begin{document}
\maketitle
\tableofcontents
\newpage

\section{Experiment Overview}
This document presents the detailed analysis for \textbf{Result 12}: Mobility Impact on Average End-to-End Delay,
as part of the ACLIF cross-layer IoT optimization framework evaluation.

\section{Result Figure}
\begin{figure}[htbp]
  \centering
  \includegraphics[width=0.92\textwidth]{../figures/fig_mobility.png}
  \caption{Mobility Impact on Average End-to-End Delay. Results obtained from NS3 v3.38 simulations over a
  100-node heterogeneous IoT network ($200\times200$~m$^2$), averaged over 10 runs
  with 95\% confidence interval shading.}
  \label{fig:mobility}
\end{figure}

\section{Analysis and Discussion}
Three mobility conditions are evaluated: static deployment, constant velocity
($v=2$~m/s), and random waypoint (RWP, $v_{\max}=5$~m/s, pause time=10~s).
ACLIF maintains the lowest delay in all three conditions. Under RWP, ACLIF
achieves 126.8~ms vs 194.8~ms for E-GLBR (34.9\% improvement). ACLIF's
staleness-aware ICSA phase discards observations older than $2T_c=1$~s,
preventing routing decisions from relying on stale position-dependent SNR estimates
under mobility, while periodic $k$-means re-clustering every 50 rounds maintains
an up-to-date relay topology.

\subsubsection*{Staleness Penalty}
\begin{equation}
\psi_i(t) = \exp\!\left(-\frac{\max(0,\; t - t_i^\mathrm{recv} - \Delta_{\max})}{\tau_s}\right)
\end{equation}

\subsubsection*{Summary: Delay Under RWP Mobility}
\begin{tabular}{lcc}
\toprule
Method & RWP Delay (ms) & vs ACLIF \\
\midrule
E-GLBR & 194.8 & +53.6\% \\
FDRL   & 167.9 & +32.4\% \\
FACSO  & 212.3 & +67.4\% \\
EEHCT  & 178.2 & +40.5\% \\
IPSO   & 162.4 & +28.1\% \\
\textbf{ACLIF} & \textbf{126.8} & --- \\
\bottomrule
\end{tabular}

\section{Experimental Configuration}
\begin{tabular}{ll}
\toprule
\textbf{Parameter} & \textbf{Value} \\
\midrule
Simulator      & NS3 v3.38 \\
Nodes          & 100 sensor nodes \\
Area           & $200 \times 200$~m$^2$ \\
PHY standard   & IEEE 802.15.4 \\
MAC protocol   & CSMA/CA \\
Packet payload & 512 bytes \\
Packet rate    & 5 pkt/s (default) \\
Runs per config & 10 (different seeds) \\
CI             & 95\% Student-$t$ \\
\bottomrule
\end{tabular}

\section{Baselines}
ACLIF is compared against: E-GLBR \cite{Benelhouri2023}, FDRL \cite{Suresh2024},
FACSO \cite{Yang2024204}, EEHCT \cite{Chaurasiya202325941}, and IPSO \cite{Zhang2025}.

\begin{thebibliography}{9}
\bibitem{Benelhouri2023} Benelhouri et al., ``E-GLBR Genetic Algorithm Routing,'' 2023.
\bibitem{Suresh2024} Suresh et al., ``FDRL Federated Deep RL Routing,'' 2024.
\bibitem{Yang2024204} Yang et al., ``FACSO Fuzzy Adaptive Cross-Layer Scheduling,'' 2024.
\bibitem{Chaurasiya202325941} Chaurasiya et al., ``EEHCT Energy-Efficient Clustering,'' 2023.
\bibitem{Zhang2025} Zhang et al., ``IPSO-Based Multi-Hop Routing,'' 2025.
\end{thebibliography}

\end{document}
